\documentclass[11pt,a4paper]{article}
\usepackage[margin=25mm]{geometry}
\usepackage{amsmath,amssymb,amsthm,hyperref}
\setlength{\emergencystretch}{2em}
\newtheorem{lemma}{Lemma}\newtheorem{theorem}[lemma]{Theorem}
\newcommand{\gen}[1]{\langle#1\rangle}
\title{The first split in an anchored string action}
\author{Human Compression project: a component of the L019 recovery}
\date{}
\begin{document}\maketitle
Let $G=\gen{\rho_0,\ldots,\rho_{r-1}}$ be a transitive
permutation group on $n$ points, generated by nonidentity
involutions satisfying the string relations. Assume every maximal
omission $G_i=\gen{\rho_j:j\ne i}$ is intransitive. An edge
of colour $i$ is eligible if its endpoints lie in different
$G_i$-orbits. An endpoint anchor $e$ is moved by $\rho_0$ and
fixed by every later mark.

We acquire the paired-cut bound, the unique-crossing lemma,
the actual head-support proof and the ambient interval deficit.
These are statements about the full permutation action; they
require neither parity nor the intersection property.

\begin{lemma}[The anchor determines one side]\label{anchor}
Suppose an endpoint anchor $e$ exists. The noncommutation
diagram is connected. Suppose a unique
eligible edge of colour $j>0$ has omission orbits $A,B$, with
$e\in A$. Then $G_{<j}$ is transitive on $A$ and $G_{>j}$
fixes $A$ pointwise. If this split is not perfect, both sides
are nontrivial and some mark of index less than $j$ acts on
both sides.
\end{lemma}
\begin{proof}
The component of the noncommutation diagram containing the
first mark generates a group $H$. Every other component
centralizes $H$ and fixes $e$. Thus the full orbit of $e$
is its $H$-orbit, so $H$ is transitive. A permutation
centralizing a transitive group and fixing one point fixes
all points. Hence there is no other nontrivial component.

Put $L=G_{<j}$ and $R=G_{>j}$. These groups commute, $R$
fixes $e$, and $LR=G_j$ is transitive on $A$. Consequently
$e^L=A$, and centralization makes $R$ fix $A$ pointwise.
The side $A$ is nontrivial since $\rho_0$ moves $e$.
If $B$ were a singleton, all of $L$ would fix it and the
split would be perfect. Otherwise nonperfection means that
some left mark acts on $B$. To find a shared mark, if no left mark
were shared then each would act on precisely one side.
Their side could not change at adjacent positions, by the
connected noncommutation diagram. The first left mark moves
$e\in A$, so all would be supported on $A$, a contradiction.
\end{proof}

\begin{lemma}[The head bound in its exact rank range]\label{head}
In the setting of the acquired head-support proof, suppose
a unique cut is at $i$, its sides are nontrivial, the group
$G_{<i}$ fixes its second side, and some later mark is shared.
If $2r>n$, the final shared index $h$ satisfies $h<r-1$.
For the support $X$ of $G_{>h}$, the same proof gives
\[
                       2h\le n-|X|-1.                 \tag{1}
\]
At $i>0$, the tail is supported on the second side. At
$i=0$, call the side containing that tail the second side.
\end{lemma}
\begin{proof}
The first five lemmas of the acquired head-support manuscript
have no global rank hypothesis. Its last-mark lemma says
$n\ge2r$ if the final mark is shared, or if $i>0$ and it
acts on the first side. Both alternatives contradict the
present hypothesis. Connectedness forces all marks after
$h$ onto the same side: two adjacent nonidentity marks with
disjoint supports would commute. At $i=0$ the side names
may be interchanged. This is exactly the placement used
in the acquired proof.

The remaining fixed-point count there is unchanged. At
$i=0$, its two disjoint path pieces have respectively
$h+1$ and $h$ fixed points. At $i>0$, its transitive-side
case gives $|A|+1\ge2i+3=2h+1$; its intransitive-side
case gives $|A|+(h-i)\ge2h+1$. These inequalities use
only the commuting-orbit, internal-pair and path lemmas,
whose full weak equality cases are retained. Thus the
tail fixes at least $2h+1$ points, proving (1).
No step of this oriented argument uses the stronger
rank assumption needed elsewhere to obtain its orientation.
\end{proof}

\begin{theorem}\label{first}
Suppose an endpoint anchor exists and $2r>n$. There is
a colour with a unique eligible edge. The least such
colour $j$ is a perfect split. If $A$ is the side containing
$e$, then
\[
                              2j\le |A|-1.             \tag{2}
\]
\end{theorem}
\begin{proof}
Full transitivity ensures at least one eligible edge of
each colour. If all had at least two, the acquired paired-cut
bound would give $2r\le n$. Thus $j$ exists.
If $j=0$, the anchor is a singleton $G_0$-orbit.
The unique-crossing lemma gives exactly one other orbit,
so the split is perfect and (2) is equality.

Let $j>0$, and suppose this split is not perfect.
Use Lemma~\ref{anchor} and reverse all labels.
The crossing index is $i'=r-1-j$; its first side is the
old $B$, and its second side is the old $A$.
The earlier group in this reversed order is the old
$G_{>j}$, which fixes $A$. Lemma~\ref{head} applies.
If $g$ is the first shared index in the original order,
the reversed final shared index is $h'=r-1-g$.
Thus $0<g<j$, and, with $X=\operatorname{supp}(G_{<g})$,
\[
                     2(r-1-g)\le n-|X|-1.             \tag{3}
\]
For $j<r-1$ the support placement is already fixed by
Lemma~\ref{head}. If $j=r-1$, its reversed crossing is
zero. The final reversed tail contains the original
$\rho_0$, which moves $e\in A$, so that tail is on $A$.
This checks the endpoint naming as well.

Every colour in the nonempty prefix $\{0,\ldots,g-1\}$
has at least two ambient eligible edges, by minimality
of $j$. Its group is intransitive, being contained in
$G_j$, and has at least one nontrivial orbit. The acquired
ambient interval deficit therefore gives
\[
                              2g\le |X|-1.             \tag{4}
\]
Adding (3) and (4) yields $2r-2\le n-2$, a contradiction.
This proves perfection.

If $j>0$, apply the same interval deficit to the full
prefix $\{0,\ldots,j-1\}$. At a perfect split it is
transitive on $A$, fixes $B$, and therefore has precisely
one nontrivial orbit and support $A$. Its $j$ colours
are all paired by minimality, giving (2).
\end{proof}

\paragraph{Attribution and scope.}
The fracture and split framework is due to the cited
Cameron--Fernandes--Leemans sources and their earlier work
with Mixer. The complete acquired orbit-cut, path and
interval proofs retain those attributions and their
separately disclosed contributions. The present deduction
uses the endpoint anchor to fix the orientation, then pays
the full interval orbit deficit at the first unique cut.
It replaces the longer proposed two-anchor density route
for an anchored parent, without needing a finite exceptional
list or a quotient intersection property.

For a critical even-degree symmetric factor of rank $m/2+1$,
and for its anchored odd-degree tail of rank $(d+1)/2$, the
theorem supplies a perfect split and (2). It does not force
that split to be interior, make its internal factor trivial,
or pay its leaf kernels. In particular $j=0$ is allowed.
The general sufficient-split and stable correspondence
obligations remain separate. No new computation, whole
original, conjecture, Lean or measurement credit is claimed.
\end{document}
